\section{Numerical Derivatives and Interpolation}
\thispagestyle{plain}

The derivative of a function $f: \mathbb{R} \to \mathbb{R}$
at $x \in \mathbb{R}$ is defined as

\begin{equation}
    f'(x) := \lim_{\delta \to 0} \frac{f(x+\delta)-f(x)}{\delta}
\end{equation}

and I assume you to be familiar with various \textit{symbolic} differentiation
rules, e.g. $\frac{d\sin}{dx}(x) = \cos(x)$.

Consider the function $f$ is the output of a (possibly complex) computer program - how
would you calculate the derivative? There are two general options:

\begin{itemize}
    \item \textbf{finite differencing}: one approximates the derivative at a given point $x$ based on evaluations of $f$
    \item \textbf{automatic differentiation}: the primitive building blocks of the program come with pre-implemented derivatives,
    the chain rule is applied to differentiate $f$
\end{itemize}

Details on automatic differentiation - the backbone of modern machine 
learning - can be found in Sec. \ref{app:autodiff} in the appendix.
Here we focus on \textbf{finite differencing}.

\bluebox{\textbf{Aim:} Approximate $f'(x)$ based on evaluations of $f$.}

Such approximations can be derived based on Taylor expansion.

\subsection{Taylor expansion}

In the vicinity of a point, Taylor expansion approximates a function with a polynomial.

Let $I \subset \mathbb{R}$ and $x_0$ an inner point of $I$. Consider a function $f: \mathbb{R} \to \mathbb{R}$
which is $n$ times continuously differentiable on $I$ and $n+1$ times differentiable inside $I$.

For all $x \in I$ we can Taylor expand

\begin{equation}
    f(x) = T_n(x) + R_n(x)
\end{equation}

where $T_n$ is the degree-$n$ polynomial

\begin{equation}
    T_n(x) = \sum_{k = 0}^{n} \frac{f^{(k)}(x_0)}{k!} (x - x_0)^k
\end{equation}

and $R_n(x)$ is a remainder, in \textit{Lagrange form} given by

\begin{equation}
    R_n(x) = \frac{f^{(n+1)}(\eta)}{(n+1)!}(x-x_0)^{n+1} \in \mathcal{O}(|x-x_0|^{n+1})
\end{equation}

where we used Landaus big-$\mathcal{O}$ notation.

\greybox{\textbf{Landau notation}:
\begin{itemize}
    \item little $\mathcal{o}$: $f \in \mathcal{o}(g) \iff \lim_{x \to a} \left| \frac{f(x)}{g(x)} \right| = 0$ ($f$ grows slower than $g$)
    \item big $\mathcal{O}$: $f \in \mathcal{O}(g) \iff \limsup_{x \to a} \left| \frac{f(x)}{g(x)} \right| < \infty$ ($f$ grows at most as fast as $g$)
\end{itemize}
}

Consider $x = x_0 + \Delta x$, then we compactly write

\begin{equation}
    \boxed{
    f(x_0 + \Delta x) = \sum_{k = 0}^{n} \frac{f^{(k)}(x_0)}{k!} \Delta x^k + \mathcal{O}(\Delta x^{n+1})
    }
\end{equation}

The interpolation error is of order $n+1$.

\subsection{Finite differencing}

Let us discretize a function $f:\mathbb{R} \to \mathbb{R}$ to a grid of evaluations $f_i = f(x_i)$ at points
$x_i$ with $x_{i+1} = x_{i} + \Delta x$. We will denote $f'_i \equiv \frac{df}{dx}(x_i), f''_i \equiv \frac{d^2f}{dx^2}(x_i),...$

\subsubsection{Derivation of finite differencing formulas from Taylor expansion}

By Taylor expansion, we can derive finite difference formulas for the derivatives.

\begin{itemize}
    \item By a forward Taylor expansion up to first order
            \begin{equation}
                f_{i+1} = f_{i} + \Delta x \, f'_i + \mathcal{O}(\Delta x^2) \quad \rightarrow \quad f'_i = \frac{f_{i+1} - f_{i}}{\Delta x} + \mathcal{O}(\Delta x)
            \end{equation}
          we yield the \textbf{first order forward difference} formula.
    \item Likewise,
        \begin{equation}
            f_{i-1} = f_{i} - \Delta x \, f'_i + \mathcal{O}(\Delta x^2) \quad \rightarrow \quad f'_i = \frac{f_{i} - f_{i-1}}{\Delta x} + \mathcal{O}(\Delta x)
        \end{equation}
          yields the \textbf{first order backward difference} formula.
\end{itemize}

By considering more evaluation points, we can reach higher orders in the \textit{truncation error}:

\begin{equation}
    \left. \begin{matrix}
        f_{i+1} = f_{i} + \Delta x \, f'_i + \frac{1}{2} f''_i \Delta x^2 + \mathcal{O}(\Delta x^3) \\
        f_{i-1} = f_{i} - \Delta x \, f'_i + \frac{1}{2} f''_i \Delta x^2 + \mathcal{O}(\Delta x^3)
    \end{matrix} \right\} \rightarrow f'_i = \frac{f_{i+1} - f_{i-1}}{2 \Delta x} + \mathcal{O}(\Delta x^2)
\end{equation}

here the \textbf{second order central difference} formula.

We can also derive approximations for \textbf{higher-order derivatives}

\begin{equation}
    \left. \begin{matrix}
        f_{i+1} = f_{i} + \Delta x \, f'_i + \frac{1}{2} f''_i \Delta x^2 + \mathcal{O}(\Delta x^3) \\
        f_{i-1} = f_{i} - \Delta x \, f'_i + \frac{1}{2} f''_i \Delta x^2 + \mathcal{O}(\Delta x^3)
    \end{matrix} \right\} \rightarrow f''_i = \frac{f_{i+1} - 2 f_{i}  + f_{i-1}}{\Delta x^2} + \mathcal{O}(\Delta x)
\end{equation}

\subsubsection{Truncation vs. floating point errors}

We would like to choose a small $\Delta x$ to reduce the truncation error. However, for very
small $\Delta x$ cancellation occurs (as we subtract very similar numbers,
e.g. $f_{i+1}$ and $f_{i}$, magnified by dividing by a small number) 
and the error in our approximation is dominated by the floating 
point error (see e.g. Fig. 3 in \cite{baydin18} and \cite{sauer25}).

Let us consider the \textbf{central difference formula} and assume
floating point representations

\begin{equation}
    \hat{f}_i = f_i + \epsilon_i, \quad \text{assume} \, |\epsilon_i| \leq \epsilon_\text{mach}
\end{equation}

such that

\begin{equation}
    f'_i = \frac{\hat{f}_{i+1} - \epsilon_{i + 1} - \hat{f}_{i-1} + \epsilon_{i - 1}}{2 \Delta x} + \mathcal{O}(\Delta x^2) = \frac{\hat{f}_{i+1} - \hat{f}_{i-1}}{2 \Delta x} + \mathcal{O}(\Delta x^2) + \frac{\epsilon_{i - 1} - \epsilon_{i + 1}}{2\Delta x}
\end{equation}

where by our assumption $\epsilon_{i - 1} - \epsilon_{i + 1} \leq 2\epsilon_\text{mach}$

\begin{equation}
    \underbrace{f'_i}_{\text{exact derivative}} = \underbrace{\frac{\hat{f}_{i+1} - \hat{f}_{i-1}}{2 \Delta x}}_{\text{finite difference}} + \underbrace{\mathcal{O}(\Delta x^2)}_{\text{truncation error}} + \underbrace{\mathcal{O}(\frac{\epsilon_\text{mach}}{\Delta x})}_{\text{floating point error}}
\end{equation}

\greybox{\textbf{Task to the reader:} Show that the error term is minimal for $\Delta x \sim \epsilon_\text{mach}^{1/3}$.}

\subsection{Discretize and solve - the power of finite differencing}

Scientific principles are typically captured in the form of differential equations, for 
instance the diffusion equation in two dimensions

\begin{equation}
    \partial_t u = K \, (\partial_x^2 + \partial_y^2) u
\end{equation}

Finite differencing is the simplest option for solving such equations. The two principal
discretization options across the independent variables are

\begin{itemize}
    \item discretize onto a grid
    \item discretize into an iteration
\end{itemize}

Spatial variables are typically discretized onto a grid and time is discretized into
an iteration. For the 2D diffusion equation above, we discretize space into a regular
$2D$ grid $u_{ij} = u(x_i,y_j)$ 

\begin{equation}
    K \, ((\partial_x^2 + \partial_y^2) u)_{ij} \approx K \, \left( \frac{u_{i+1,j} - 2 u_{ij} + u_{i-1,j}}{\Delta x^2} + \frac{u_{i,j+1} - 2 u_{ij} + u_{i,j-1}}{\Delta y^2} \right)
\end{equation}

and time into an explicit iteration

\begin{equation}
    \partial_t u_{ij}^{(n)} \approx \frac{u_{ij}^{(n+1)} - u_{ij}^{(n)}}{\Delta t} \approx K \, ((\partial_x^2 + \partial_y^2) u)_{ij}^{(n)} \quad \rightarrow \quad u_{ij}^{(n+1)} \approx u_{ij}^{(n)} + \Delta t \, K \, ((\partial_x^2 + \partial_y^2) u)_{ij}^{(n)}
\end{equation}

so together with our spatial finite difference approximation, we obtain the explicit finite difference scheme

\begin{equation}
    u_{ij}^{(n+1)} \approx u_{ij}^{(n)} + \Delta t \, K \, \left( \frac{u^{(n)}_{i+1,j} - 2 u^{(n)}_{ij} + u^{(n)}_{i-1,j}}{\Delta x^2} + \frac{u^{(n)}_{i,j+1} - 2 u^{(n)}_{ij} + u^{(n)}_{i,j-1}}{\Delta y^2} \right)
\end{equation}

which is started from some initial conditions $u_{ij}^{(0)} = u_{ij}(t = t_0)$.

\subsection{Interpolation}
Consider we have data points $f_i = f(x_i), i = 0,..,m$ and would like to 
find a polynomial $p \in \Pi_m$ passing through those points $p(x_i) = f(x_i)$.

\bluebox{\textbf{Theorem of Polynomial Interpolation:} Let $(x_0, f_0),... , (x_m, f_m)$ be n points in the
plane with distinct $x_i$. Then there exists one and only one polynomial $p$ of degree $m$ or
less that satisfies $p(x_i) = f_i$ for $i = 0,...,m$. See \cite[Theorem 2.3]{sauer25}.}

In fluid simulations, for instance,
we might deal with cell-centered fluid quantities and would like to interpolate to the
cell boundaries to calculate intercell fluxes.

\subsubsection{Evaluating polynomials efficiently - Horner scheme}

Let us start with the efficient evaluation of polynomials.

Consider $p \in \Pi_m$, so

\begin{equation}
    p(x) = b_0 + b_1 x + b_2 x^2 + ... + b_m x^m
\end{equation}

We could obtain all the powers $(x^2,...,x^m)$ by $m-1$ multiplications
and could then evaluate $p(x)$ with another $m$ multiplications
of the powers with the coefficients, in total $2m - 1$ multiplications
(and $m$ computationally cheaper additions).

We can rewrite polynomials to reduce the number of multiplications

\begin{itemize}
    \item $p(x) = b_0 + b_1 x + b_2 x^2 = b_0 + (b_1 + b_2 x) x$ ($m = 2$, $2$ multiplications, $2$ additions)
    \item $p(x) = b_0 + b_1 x + b_2 x^2 + b_3 x^3 = b_0 + (b_1 + (b_2 + b_3 x) x) x$ ($m = 3$, $3$ multiplications, $3$ additions)
    \item $\vdots$
    \item $p(x) = (\dots (b_m x + b_{m-1})x + \dots )$ ($m$ multiplications, $m$ additions)
\end{itemize}

- only $m$ multiplication are necessary in the \textbf{Horner scheme}.

The Horner scheme can be implemented in the recurrent form

\begin{equation}
    z_j = b_j + x z_{j+1}, z_m = b_m \quad \rightarrow \quad z_0 = p(x)
\end{equation}

\subsubsection{Langrange interpolation}

The simplest interpolation scheme is obtained based on Lagrange
polynomials

\begin{equation}
    l_j(x) = \prod_{k = 0, k \neq j}^{m} \frac{x-x_k}{x_j - x_k} \quad \rightarrow \quad l_j(x_i) = \delta_{ij} = \begin{cases}
        1, &i = j, \\
        0, &i \neq j
    \end{cases}
\end{equation}

designed such that $l_j$ is $1$ on $x_j$ but zero on all other supporting points.

The interpolation becomes trivial

\begin{equation}
    p(x) = f(x_0) l_0(x) + f(x_1) l_1(x) + \dots + f(x_m) l_m(x)
\end{equation}

\redbox{\textbf{Problems:} 
\begin{itemize}
    \item The computational cost of evaluating the Lagrange polynomial is $\mathcal{O}(m^2)$.
    \item Cancellation problems in the evaluation of $l_j$.
\end{itemize}}

\subsubsection{Newton interpolation}



\subsubsection{Outlook}

\begin{itemize}
    \item splines
    \item neural networks - not a polynomial basis
    \item regression tasks and optimization
\end{itemize}

\pagebreak